\documentclass[12pt, letterpaper]{article}

\usepackage[utf8]{inputenc}
\usepackage[T1]{fontenc}
\usepackage[spanish]{babel}
\usepackage{geometry}
\usepackage{amsmath, amssymb, amsthm}
\usepackage{booktabs}
\usepackage{array}
\usepackage{microtype}
\usepackage{setspace}
\usepackage{parskip}
\usepackage{titlesec}
\usepackage{fancyhdr}
\usepackage{enumitem}
\usepackage{bm}
\usepackage{mathtools}
\usepackage{tcolorbox}
\usepackage{hyperref}

\tcbuselibrary{skins, breakable}

\geometry{top=2.8cm, bottom=3.0cm, left=3.2cm, right=3.2cm}
\setstretch{1.20}

\pagestyle{fancy}
\fancyhf{}
\fancyhead[C]{\small\textsc{Econometria · Gujarati --- Ejercicio 11.3}}
\fancyfoot[C]{\small\thepage}
\renewcommand{\headrulewidth}{0.4pt}
\renewcommand{\footrulewidth}{0pt}

\titleformat{\section}
  {\large\bfseries\scshape}{\thesection.}{0.6em}{}
  [\vspace{-4pt}\rule{\linewidth}{0.4pt}]
\titleformat{\subsection}{\normalsize\bfseries}{\thesubsection.}{0.5em}{}
\titleformat{\subsubsection}{\normalsize\itshape}{}{0em}{}

\newtheorem{prop}{Proposicion}
\theoremstyle{remark}
\newtheorem{obs}{Observacion}

\newtcolorbox{clave}{
  colback=white, colframe=black,
  boxrule=0.5pt, arc=3pt,
  left=8pt, right=8pt, top=6pt, bottom=6pt,
  breakable
}

\newcommand{\E}{\mathbb{E}}
\newcommand{\Var}{\operatorname{Var}}
\newcommand{\plim}{\operatorname{plim}}

\begin{document}

\begin{center}
  {\LARGE\bfseries Modelos No Lineales en los Parametros:}\\[6pt]
  {\LARGE\bfseries MCO, Minimos Cuadrados No Lineales}\\[6pt]
  {\LARGE\bfseries y Metodos de Estimacion Alternativos}\\[12pt]
  {\large\itshape Por que MCO falla y como estimar modelos con raices}\\[14pt]
  \rule{0.55\linewidth}{0.6pt}\\[8pt]
  {\normalsize Ejercicio~11.3 --- \textit{Econometria}, Gujarati \& Porter, 5.a ed.}\\[4pt]
  {\small\textsc{Con derivacion de las ecuaciones normales no lineales, algoritmo de Gauss-Newton
  y comparacion de metodos}}
\end{center}

\vspace{10pt}

\noindent\textbf{Planteamiento.}\quad
Se proponen dos modelos donde la variable dependiente es el valor absoluto de
los residuos MCO ($|\hat u_i|$) como proxy de la desviacion estandar condicional:

\begin{align}
  |\hat u_i| &= \sqrt{\beta_1 + \beta_2 X_i} + v_i
  \tag{I}\label{eq:modelo1}\\[4pt]
  |\hat u_i| &= \sqrt{\beta_1 + \beta_2 X_i^2} + v_i
  \tag{II}\label{eq:modelo2}
\end{align}

\noindent Se pregunta si se pueden estimar los parametros $\beta_1$ y $\beta_2$
por MCO y, si no, que metodos alternativos existen.

\section{Inciso (a): ¿puede MCO estimar estos modelos?}

\subsection{Respuesta}

\textbf{No}. El metodo de MCO no puede aplicarse directamente a
los modelos~(\ref{eq:modelo1}) y~(\ref{eq:modelo2}) porque ambos son
\textbf{no lineales en los parametros}.

\subsection{Condicion de linealidad en los parametros}

MCO minimiza la suma de cuadrados residuales:
\begin{equation}
  S(\boldsymbol\beta) = \sum_{i=1}^n v_i^2 = \sum_{i=1}^n [Y_i - f(X_i; \boldsymbol\beta)]^2
  \label{eq:SCR}
\end{equation}

Para que MCO produzca ecuaciones normales \emph{lineales} en $\boldsymbol\beta$
(y por tanto tenga solucion analitica explicita), la funcion $f$ debe ser
lineal en los parametros:
\begin{equation}
  f(X_i; \boldsymbol\beta) = \beta_1 g_1(X_i) + \beta_2 g_2(X_i) + \cdots
  \label{eq:lineal_param}
\end{equation}
para algunas funciones $g_j$ de los datos. En este caso las derivadas parciales:
\begin{equation}
  \frac{\partial f}{\partial \beta_j} = g_j(X_i)
  \label{eq:deriv_lineal}
\end{equation}
no dependen de los parametros, y las ecuaciones normales
$\partial S/\partial \beta_j = 0$ forman un sistema lineal resoluble
analitica­mente.

\subsection{Por que los modelos (I) y (II) violan esta condicion}

En el modelo~( \ref{eq:modelo1} ), la funcion de regresion es:
\begin{equation}
  f(X_i; \beta_1, \beta_2) = \sqrt{\beta_1 + \beta_2 X_i}
  = (\beta_1 + \beta_2 X_i)^{1/2}
  \label{eq:f_modelo1}
\end{equation}

Las derivadas parciales respecto a los parametros son:
\begin{align}
  \frac{\partial f}{\partial \beta_1}
    &= \frac{1}{2}(\beta_1 + \beta_2 X_i)^{-1/2}
    = \frac{1}{2\sqrt{\beta_1 + \beta_2 X_i}}
  \label{eq:df_db1}\\[6pt]
  \frac{\partial f}{\partial \beta_2}
    &= \frac{X_i}{2}(\beta_1 + \beta_2 X_i)^{-1/2}
    = \frac{X_i}{2\sqrt{\beta_1 + \beta_2 X_i}}
  \label{eq:df_db2}
\end{align}

Ambas derivadas dependen de los propios parametros $\beta_1$ y $\beta_2$, lo que
hace que las ecuaciones normales:
\begin{equation}
  \frac{\partial S}{\partial \beta_j}
    = -2\sum_i v_i \frac{\partial f}{\partial \beta_j}
    = -2\sum_i [|\hat u_i| - \sqrt{\beta_1 + \beta_2 X_i}]
      \cdot \frac{\partial f}{\partial \beta_j}
    = 0
  \label{eq:ec_normales}
\end{equation}
sean \textbf{no lineales} en $\beta_1$ y $\beta_2$. No hay solucion analitica
cerrada; MCO ordinario no puede resolver este sistema.

\subsection{Sustitucion cuadratica: ¿resuelve el problema?}

Una tentacion es elevar al cuadrado ambos lados de~(\ref{eq:modelo1}):
\begin{equation}
  |\hat u_i|^2 = \hat u_i^2 = \beta_1 + \beta_2 X_i + w_i
  \label{eq:cuadrado}
\end{equation}
donde $w_i = v_i^2 + 2v_i\sqrt{\beta_1 + \beta_2 X_i}$. El lado izquierdo
$\hat u_i^2$ es observable y el modelo en~(\ref{eq:cuadrado}) es \emph{lineal}
en $\beta_1$ y $\beta_2$. Sin embargo, el nuevo termino de error $w_i$ es
heterosc­edastico y correlacionado con los regresores, por lo que MCO sobre
la version cuadrada no produce estimadores MELI. Es un atajo que introduce
problemas de estimacion.

\begin{clave}
  MCO no puede estimar los modelos~(\ref{eq:modelo1}) y~(\ref{eq:modelo2}) porque
  son no lineales en los parametros: las ecuaciones normales resultantes son
  no lineales en $\beta_1$ y $\beta_2$ y no tienen solucion analitica cerrada.
\end{clave}

\section{Inciso (b): metodos de estimacion alternativos}

\subsection{Metodo formal 1: Minimos Cuadrados No Lineales (MCNL)}

El metodo de \textbf{Minimos Cuadrados No Lineales} (NLLS, \textit{Nonlinear
Least Squares}) generaliza MCO al caso de modelos no lineales en los parametros.
El problema de minimizacion es el mismo:
\begin{equation}
  \hat{\boldsymbol\beta}^{\text{MCNL}}
    = \arg\min_{\boldsymbol\beta} S(\boldsymbol\beta)
    = \arg\min_{\boldsymbol\beta} \sum_i [|\hat u_i| - f(X_i; \boldsymbol\beta)]^2
  \label{eq:MCNL}
\end{equation}
pero ahora $f$ es no lineal en $\boldsymbol\beta$. La diferencia es que la
solucion se obtiene \emph{iterativamente} en lugar de analiticamente.

\subsubsection{Propiedades asintóticas del estimador MCNL}

Bajo condiciones de regularidad, el estimador MCNL es:
\begin{itemize}[leftmargin=1.6em, itemsep=3pt]
  \item \textbf{Consistente}: $\hat{\boldsymbol\beta}^{\text{MCNL}} \xrightarrow{p} \boldsymbol\beta_0$
  \item \textbf{Asintóticamente normal}: $\sqrt{n}(\hat{\boldsymbol\beta}^{\text{MCNL}} - \boldsymbol\beta_0)
    \xrightarrow{d} \mathcal{N}(\mathbf{0}, \mathbf{V})$
  \item \textbf{Asintóticamente eficiente} dentro de la clase de estimadores
    de momentos bajo homoscedasticidad.
\end{itemize}

La matriz de varianzas asintótica es:
\begin{equation}
  \mathbf{V} = \sigma^2 \left(\mathbf{F}'\mathbf{F}\right)^{-1}
  \label{eq:var_MCNL}
\end{equation}
donde $\mathbf{F}$ es la matriz de derivadas parciales evaluadas en
$\hat{\boldsymbol\beta}^{\text{MCNL}}$:
\begin{equation}
  F_{ij} = \left.\frac{\partial f(X_i; \boldsymbol\beta)}{\partial \beta_j}\right|_{\boldsymbol\beta = \hat{\boldsymbol\beta}}
  \label{eq:F_matrix}
\end{equation}

\subsection{Metodo formal 2: algoritmo de Gauss-Newton}

El algoritmo de \textbf{Gauss-Newton} es el procedimiento iterativo mas comun
para MCNL. La idea es linealizar $f$ en cada iteracion mediante una expansion
de Taylor de primer orden alrededor del valor actual $\boldsymbol\beta^{(k)}$:

\begin{equation}
  f(X_i; \boldsymbol\beta)
    \approx f(X_i; \boldsymbol\beta^{(k)})
    + \sum_j \left.\frac{\partial f}{\partial \beta_j}\right|_{\boldsymbol\beta^{(k)}}
      (\beta_j - \beta_j^{(k)})
  \label{eq:taylor}
\end{equation}

Definiendo:
\begin{align}
  \tilde Y_i^{(k)} &= |\hat u_i| - f(X_i; \boldsymbol\beta^{(k)})
    + \sum_j F_{ij}^{(k)} \beta_j^{(k)}
  \label{eq:Ytilde}\\
  \tilde X_{ij}^{(k)} &= F_{ij}^{(k)} = \left.\frac{\partial f}{\partial \beta_j}\right|_{\boldsymbol\beta^{(k)}}
  \label{eq:Xtilde}
\end{align}

La actualizacion en la iteracion $k+1$ se obtiene por MCO de
$\tilde Y_i^{(k)}$ sobre $\tilde X_{ij}^{(k)}$:

\begin{equation}
  \boldsymbol\beta^{(k+1)}
    = \boldsymbol\beta^{(k)}
    + \left[(\tilde{\mathbf{X}}^{(k)})'\tilde{\mathbf{X}}^{(k)}\right]^{-1}
      (\tilde{\mathbf{X}}^{(k)})' \tilde{\mathbf{Y}}^{(k)}
  \label{eq:GN_update}
\end{equation}

El algoritmo se repite hasta convergencia: $\|\boldsymbol\beta^{(k+1)} - \boldsymbol\beta^{(k)}\| < \varepsilon$
para algun criterio de tolerancia $\varepsilon$.

\subsubsection{Aplicacion al modelo (I)}

Para el modelo~(\ref{eq:modelo1}), las derivadas~(\ref{eq:df_db1}) y~(\ref{eq:df_db2})
evaludas en $\boldsymbol\beta^{(k)}$ son:
\begin{align}
  \tilde X_{i1}^{(k)} &= \frac{1}{2\sqrt{\beta_1^{(k)} + \beta_2^{(k)} X_i}}
  \label{eq:GN_X1}\\
  \tilde X_{i2}^{(k)} &= \frac{X_i}{2\sqrt{\beta_1^{(k)} + \beta_2^{(k)} X_i}}
  \label{eq:GN_X2}
\end{align}

Estas son funciones de los datos observados y de los valores actuales de los
parametros. En cada iteracion se recalculan y se aplica MCO sobre las
variables pseudo-lineales $(\tilde Y^{(k)}, \tilde X_1^{(k)}, \tilde X_2^{(k)})$.

\subsection{Metodo formal 3: Newton-Raphson y Marquardt}

El algoritmo de \textbf{Newton-Raphson} usa la matriz Hessiana de $S(\boldsymbol\beta)$
para la actualizacion:
\begin{equation}
  \boldsymbol\beta^{(k+1)}
    = \boldsymbol\beta^{(k)}
    - \left[\mathbf{H}^{(k)}\right]^{-1} \mathbf{g}^{(k)}
  \label{eq:NR}
\end{equation}
donde $\mathbf{g}^{(k)} = \nabla_{\boldsymbol\beta} S(\boldsymbol\beta^{(k)})$
es el gradiente y $\mathbf{H}^{(k)} = \nabla^2_{\boldsymbol\beta} S(\boldsymbol\beta^{(k)})$
es la Hessiana. Es mas rapido que Gauss-Newton cerca del maximo pero requiere
calcular segundas derivadas.

El algoritmo de \textbf{Marquardt} (o Levenberg-Marquardt) combina
Gauss-Newton y descenso de gradiente:
\begin{equation}
  \boldsymbol\beta^{(k+1)}
    = \boldsymbol\beta^{(k)}
    + \left[(\tilde{\mathbf{X}}^{(k)})'\tilde{\mathbf{X}}^{(k)}
            + \mu^{(k)}\mathbf{I}\right]^{-1}
      (\tilde{\mathbf{X}}^{(k)})' \tilde{\mathbf{Y}}^{(k)}
  \label{eq:Marquardt}
\end{equation}
donde $\mu^{(k)} \geq 0$ es un parametro de amortiguacion. Para $\mu^{(k)} = 0$
se recupera Gauss-Newton; para $\mu^{(k)} \to \infty$ se obtiene descenso de
gradiente. Este metodo es mas robusto que Gauss-Newton puro y es el estandar
en la mayoria de los paquetes de software.

\subsection{Metodo formal 4: Maxima Verosimilitud (MV)}

Bajo el supuesto de que $v_i \sim \mathcal{N}(0, \sigma^2)$, la funcion de
log-verosimilitud para el modelo~(\ref{eq:modelo1}) es:
\begin{equation}
  \ell(\boldsymbol\beta, \sigma^2)
    = -\frac{n}{2}\ln(2\pi) - \frac{n}{2}\ln\sigma^2
    - \frac{1}{2\sigma^2}\sum_i [|\hat u_i| - \sqrt{\beta_1 + \beta_2 X_i}]^2
  \label{eq:loglik}
\end{equation}

Maximizar~(\ref{eq:loglik}) respecto a $\boldsymbol\beta$ es equivalente a
minimizar $S(\boldsymbol\beta)$ en~(\ref{eq:MCNL}): los estimadores MV y MCNL
coinciden bajo normalidad. Los estimadores MV tienen propiedades asintóticas
optimas (consistencia, normalidad asintótica, eficiencia asintótica).

\subsection{Metodo informal: busqueda en grilla (prueba y error)}

Un metodo informal es la \textbf{busqueda en grilla} (\textit{grid search}):

\begin{enumerate}[leftmargin=1.8em, itemsep=4pt]
  \item Definir una rejilla de valores candidatos para $(\beta_1, \beta_2)$, por
    ejemplo $\beta_1 \in \{0.1, 0.2, \ldots, 2.0\}$ y $\beta_2 \in \{0.01, 0.02, \ldots, 0.5\}$.
  \item Para cada par $(\beta_1, \beta_2)$ de la rejilla, calcular:
    \begin{equation}
      S(\beta_1, \beta_2) = \sum_i [|\hat u_i| - \sqrt{\beta_1 + \beta_2 X_i}]^2
    \end{equation}
  \item Seleccionar el par que minimiza $S(\beta_1, \beta_2)$.
  \item Refinar la busqueda alrededor del minimo encontrado con una grilla
    mas fina.
\end{enumerate}

Este metodo es computacionalmente intensivo pero conceptualmente simple. No
garantiza encontrar el minimo global si la funcion $S$ tiene multiples minimos
locales, y es impractico en alta dimension.

\subsection{Comparacion de metodos}

\begin{table}[ht]
\centering
\caption{Comparacion de metodos para estimar modelos no lineales en los parametros}
\label{tab:metodos}
\small
\begin{tabular}{p{3.2cm} p{3.5cm} p{3.5cm} p{3cm}}
\toprule
\textbf{Metodo} & \textbf{Ventajas} & \textbf{Desventajas} & \textbf{Cuando usar} \\
\midrule
MCNL / Gauss-Newton
  & Solucion sistematica; propiedades asintóticas conocidas
  & Requiere valores iniciales; puede no converger
  & Estandar para modelos no lineales \\[4pt]
Marquardt
  & Mas robusto que G-N; converge en mas casos
  & Ligeramente mas costoso computacionalmente
  & Cuando G-N falla \\[4pt]
Maxima Verosimilitud
  & Optimo asintóticamente bajo normalidad
  & Requiere supuesto distribucional; mas complejo
  & Cuando se conoce la distribucion de $v_i$ \\[4pt]
Busqueda en grilla
  & Simple; no requiere derivadas; encuentra mínimo global si la grilla es fina
  & Lenta; imprecisa; no escala a muchos parametros
  & Diagnostico inicial; pocos parametros \\
\bottomrule
\end{tabular}
\end{table}

\begin{clave}
  El metodo recomendado es MCNL implementado mediante el algoritmo de
  Marquardt, disponible en la mayoria de los paquetes econometricos (Stata:
  \texttt{nl}; R: \texttt{nls}; Python: \texttt{scipy.optimize.curve\_fit}).
  Se necesitan valores iniciales razonables para los parametros; una busqueda
  en grilla puede servir como punto de partida.
\end{clave}

\appendix

\section{Derivacion de las ecuaciones normales del MCNL para el modelo (I)}

Para el modelo~(\ref{eq:modelo1}), las condiciones de primer orden
$\partial S/\partial \beta_j = 0$ son:

\begin{align}
  \frac{\partial S}{\partial \beta_1}
    &= -2\sum_i [|\hat u_i| - (\beta_1 + \beta_2 X_i)^{1/2}]
       \cdot \frac{1}{2(\beta_1 + \beta_2 X_i)^{1/2}}
    = 0
  \label{eq:FOC1}\\[6pt]
  \frac{\partial S}{\partial \beta_2}
    &= -2\sum_i [|\hat u_i| - (\beta_1 + \beta_2 X_i)^{1/2}]
       \cdot \frac{X_i}{2(\beta_1 + \beta_2 X_i)^{1/2}}
    = 0
  \label{eq:FOC2}
\end{align}

Simplificando:
\begin{align}
  \sum_i \frac{|\hat u_i|}{(\beta_1 + \beta_2 X_i)^{1/2}}
    &= \sum_i 1 = n
  \label{eq:FOC1_simp}\\[6pt]
  \sum_i \frac{X_i |\hat u_i|}{(\beta_1 + \beta_2 X_i)^{1/2}}
    &= \sum_i X_i
  \label{eq:FOC2_simp}
\end{align}

Las ecuaciones~(\ref{eq:FOC1_simp}) y~(\ref{eq:FOC2_simp}) son claramente no
lineales en $\beta_1$ y $\beta_2$ (los parametros aparecen dentro de la raiz
cuadrada en el denominador). No hay expresion algebraica cerrada para su
solucion: confirma que MCO ordinario no es aplicable.

\section{Valores iniciales para el algoritmo iterativo}

Una estrategia practica para obtener valores iniciales $(\beta_1^{(0)}, \beta_2^{(0)})$:

\begin{enumerate}[leftmargin=1.8em, itemsep=4pt]
  \item Elevar al cuadrado el modelo~(\ref{eq:modelo1}): $\hat u_i^2 \approx \beta_1 + \beta_2 X_i$
    (ignorando el termino de error).
  \item Estimar esta regresion lineal por MCO.
  \item Usar los estimadores MCO como valores iniciales para el algoritmo MCNL.
\end{enumerate}

Este procedimiento de dos pasos proporciona valores iniciales razonables sin
garantizar eficiencia, pero acelera la convergencia del algoritmo iterativo.

\vspace{14pt}
\noindent\rule{\linewidth}{0.4pt}

\noindent{\small\textit{Nota.} Los modelos~(\ref{eq:modelo1}) y~(\ref{eq:modelo2})
aparecen en el contexto del diagnostico de heteroscedasticidad: se postulan como
modelos para la desviacion estandar condicional $\sigma_i = \E[|u_i|]$ en funcion
de $X_i$. Si se prefiere modelar la varianza en lugar de la desviacion, el modelo
cuadratico $\hat u_i^2 = \beta_1 + \beta_2 X_i + w_i$ es lineal en los parametros
y puede estimarse directamente por MCO, aunque $w_i$ no es homoscedastico. Esta
es la base de la prueba de Breusch-Pagan (1979).}

\end{document}
